\documentclass[11pt]{article}

\usepackage[margin=1in]{geometry}
\usepackage{amsmath,amssymb,amsthm,mathtools,bm}
\usepackage{microtype}
\usepackage{hyperref}

\hypersetup{colorlinks=true,linkcolor=blue,citecolor=blue,urlcolor=blue}

\newtheorem{assumption}{Assumption}
\newtheorem{theorem}{Theorem}
\newtheorem{corollary}{Corollary}
\newtheorem{proposition}{Proposition}
\newtheorem{lemma}{Lemma}
\newtheorem{definition}{Definition}
\newtheorem{remark}{Remark}

\newcommand{\E}{\mathbb{E}}
\newcommand{\Pbb}{\mathbb{P}}
\newcommand{\R}{\mathbb{R}}
\newcommand{\cL}{\mathcal{L}}
\newcommand{\cC}{\mathcal{C}}
\newcommand{\cG}{\mathcal{G}}
\newcommand{\Normal}{\mathcal{N}}
\newcommand{\TPR}{\operatorname{TPR}}
\newcommand{\FPR}{\operatorname{FPR}}
\newcommand{\Corr}{\operatorname{Corr}}
\newcommand{\argmax}{\operatorname*{arg\,max}}
\newcommand{\argmin}{\operatorname*{arg\,min}}
\DeclareMathOperator{\sech}{sech}

\title{Theoretical Justification}
\author{}
\date{}

\begin{document}
\maketitle

Signal-agnostic model validation relies exclusively on background data. Since no unsupervised method can guarantee optimal future anomaly detection for arbitrary anomalies, a necessary assumption is that future anomalies must be deviations along a latent direction that is learnable from typical data. This assumption can be formalized as a monotone likelihood-ratio condition.

\begin{assumption}[Latent likelihood-ratio recoverability]
\label{ass:lrr}
Let \(P_0\) denote the typical distribution and \(P_1\) a future anomaly distribution. There exists a scalar latent coordinate
\[
    Z=T(X)
\]
and a nondecreasing function \(h\) such that
\[
    \frac{dP_1}{dP_0}(x)=h(T(x)).
\]
Thus larger values of \(Z\) are more anomalous in the Neyman--Pearson sense. In addition, the two typical validation domains are treated as noisy samples of the same latent structure: conditioned on \(Z\), they differ only in the randomness entailed by the sampling process.
\end{assumption}

Assumption~\ref{ass:lrr} formalizes the standard anomaly-detection premise: a good model learns a coordinate of typicality, and anomalies are detected as departures along that coordinate. Then, the goal of a model validation criterion is not to discover arbitrary unseen anomalies, but to select, among trained candidate detectors, the score that recovers this latent anomaly-relevant coordinate most reliably.

We will show that CAP selects models that best recover this coordinate because of its ability to navigate the bias-variance tradeoff that this assumption entails. In contrast to other unsupervised stability criteria, CAP favors the most informative anomaly assignments that remain stable under sampling randomness in typical data.

\section{CAP as a reliability selector for learned scores}

In our setting, trained anomaly detectors already map inputs to scalar scores.
We therefore formulate a closed-form version of Assumption~\ref{ass:lrr}
directly at the score level, modelling each detector as a noisy measurement of
a latent anomaly-relevant coordinate.

\begin{assumption}[Gaussian latent signal--noise channel]
\label{ass:channel}
Let \(\widetilde S_\phi\) denote an anomaly score produced
by a candidate detector \(\phi\), assuming that larger values are
more anomalous. Under typical data, the two validation samples satisfy
\[
    \widetilde S_\phi^{(d)}
    =
    b_\phi
    +
    a_\phi Z
    +
    \sigma_\phi \varepsilon_\phi^{(d)},
    \qquad d\in\{1,2\},
\]
where
\[
    Z,\varepsilon_\phi^{(1)},\varepsilon_\phi^{(2)}
    \stackrel{\mathrm{ind}}{\sim}\Normal(0,1),
    \qquad
    a_\phi\ge 0,
    \qquad
    \sigma_\phi\ge 0,
    \qquad
    a_\phi^2+\sigma_\phi^2>0 .
\]

The coefficient \(a_\phi\) measures the sensitivity of the score to the anomaly-relevant coordinate \(Z\), while \(\sigma_\phi \varepsilon_\phi^{(d)}\) encapsulates the remaining statistical variation that is not predictive of the anomaly direction. The standardized scores are:

\[
    S_\phi^{(d)}
    =
    \frac{\widetilde S_\phi^{(d)}-b_\phi}
    {\sqrt{a_\phi^2+\sigma_\phi^2}} .
\]
Hence
\[
    S_\phi^{(d)}
    =
    \sqrt{\rho_\phi}\,Z
    +
    \sqrt{1-\rho_\phi}\,\varepsilon_\phi^{(d)},
    \qquad d\in\{1,2\},
\]
where \(\rho_\phi\) is the fraction of typical score variance explained by \(Z\).
\[
    \rho_\phi
    =
    \frac{a_\phi^2}{a_\phi^2+\sigma_\phi^2} = \Corr\!\left(S_\phi^{(1)},S_\phi^{(2)}\right)
    \in[0,1].
\]

As per our initial assumption, future anomalies are positive shifts of the same latent coordinate:
\[
    Z\mid Y=1\sim\Normal(\delta,1),
    \qquad
    \delta>0,
\]
while the nuisance variables \(\varepsilon_\phi^{(d)}\) retain their typical-data
law. Consequently, after the same typical-data normalisation,
\[
    S_\phi\mid Y=0\sim\Normal(0,1),
    \qquad
    S_\phi\mid Y=1\sim\Normal(\delta\sqrt{\rho_\phi},1).
\]
\end{assumption}

The parameter \(\rho_\phi\) is the reliability of the learned anomaly score.
A detector has large \(\rho_\phi\) only when a large fraction of its typical
score variance is explained by the anomaly-relevant coordinate \(Z\). A detector
with \(a_\phi=0\) can have a perfectly stable marginal null histogram after
normalization, but it is uninformative for future anomalies and has
\(\rho_\phi=0\). Therefore, \(\rho_\phi\) formalises the part of score
variation that is both reproducible across validation views and aligned with
the anomaly direction.


We will show that CAP's objective is monotone with respect to this parameter, whereas marginal stability criteria can be insensitive to it. For a binary Gibbs posterior, let
\[
    q_\beta(s)=\sigma(\beta s)=\frac{1}{1+e^{-\beta s}},
    \qquad
    m_\beta(s)=2q_\beta(s)-1=\tanh\!\left(\frac{\beta s}{2}\right).
\]
The approximation capacity kernel for a pair \((s_1,s_2)\) is
\[
    \log\{q_\beta(s_1)q_\beta(s_2)+[1-q_\beta(s_1)][1-q_\beta(s_2)]\}
    =
    -\log 2+
    \log\{1+m_\beta(s_1)m_\beta(s_2)\}.
\]
Since the constant \(-\log 2\) is irrelevant for model selection, let \(L_\phi\) be the optimization objective
\[
    L_\phi(\beta)
    =
    \E\log\{1+m_\beta(S_\phi^{(1)})m_\beta(S_\phi^{(2)})\}
\]
and \(L_\phi^\star\) denote the CAP lift:
\[
    L_\phi^\star
    =
    \sup_{\beta\in[0,B]} L_\phi(\beta),
    \qquad B>0.
\]

\begin{lemma}[CAP is monotone in paired Gaussian reliability]
\label{lem:price}
Let \((U,V)\) be standard bivariate Gaussian with correlation \(\rho\in[0,1)\). For every \(\beta>0\),
\[
    C_\beta(\rho)
    =
    \E_\rho\log\{1+m_\beta(U)m_\beta(V)\}
\]
is strictly increasing in \(\rho\). Consequently, \(\sup_{\beta\in[0,B]} C_\beta(\rho)\) is strictly increasing in \(\rho\) on \([0,1]\).
\end{lemma}

\begin{proof}
Let
\[
    H_\beta(u,v)=\log\{1+m_\beta(u)m_\beta(v)\}.
\]
Price's theorem gives
\[
    \frac{d}{d\rho}\E_\rho H_\beta(U,V)
    =
    \E_\rho\left[\frac{\partial^2 H_\beta}{\partial u\,\partial v}(U,V)\right].
\]
Since \(m_\beta'(t)=(\beta/2)\sech^2(\beta t/2)>0\),
\[
    \frac{\partial^2 H_\beta}{\partial u\,\partial v}(u,v)
    =
    \frac{m_\beta'(u)m_\beta'(v)}{
        \{1+m_\beta(u)m_\beta(v)\}^2
    }
    >0.
\]
Therefore \(C_\beta(\rho)\) is strictly increasing for every fixed \(\beta>0\). The optimized objective is increasing as a supremum of increasing functions. It is strict because for \(\rho>0\), a small-\(\beta\) expansion gives
\[
    C_\beta(\rho)=\frac{\beta^2}{4}\rho+O(\beta^4),
\]
so the maximizing \(\beta\) is nonzero whenever \(\rho>0\); applying the fixed-\(\beta\) strict monotonicity at such a maximizer gives strict increase. The case \(\rho=1\) follows by continuity.
\end{proof}

\begin{theorem}[Representation-conditional CAP optimality]
\label{thm:score}
Under Assumption~\ref{ass:channel}, for every false-positive rate \(\alpha\in(0,1)\),
\[
    \TPR_\phi(\alpha)
    =
    \Phi\!\left(
        \delta\sqrt{\rho_\phi}-\Phi^{-1}(1-\alpha)
    \right).
\]
Hence \(\TPR_\phi(\alpha)\) is strictly increasing in \(\rho_\phi\). Moreover, the optimized CAP lift \(L_\phi^\star\) is strictly increasing in the same \(\rho_\phi\). Therefore, for any candidate family \(\Phi\),
\[
    \argmax_{\phi\in\Phi} L_\phi^\star
    =
    \argmax_{\phi\in\Phi} \rho_\phi
    =
    \argmax_{\phi\in\Phi} \TPR_\phi(\alpha),
\]
with equality of maximizer sets. In particular, a CAP-selected score has maximal TPR at every fixed FPR among the candidate scores satisfying the model.
\end{theorem}

\begin{proof}
Under typical data, \(S_\phi\sim\Normal(0,1)\). The FPR-\(\alpha\) threshold is therefore
\[
    \tau_\alpha=\Phi^{-1}(1-\alpha).
\]
Under anomalies,
\[
    S_\phi\mid Y=1
    =
    \sqrt{\rho_\phi}\,Z_1 + \sqrt{1-\rho_\phi}\,\varepsilon,
    \qquad Z_1\sim\Normal(\delta,1),
\]
so
\[
    S_\phi\mid Y=1\sim\Normal(\delta\sqrt{\rho_\phi},1).
\]
Thus
\[
\begin{aligned}
    \TPR_\phi(\alpha)
    &=\Pbb(S_\phi\ge \tau_\alpha\mid Y=1)  \\
    &=1-\Phi\!\left(\Phi^{-1}(1-\alpha)-\delta\sqrt{\rho_\phi}\right) \\
    &=\Phi\!\left(\delta\sqrt{\rho_\phi}-\Phi^{-1}(1-\alpha)\right).
\end{aligned}
\]
This is strictly increasing in \(\rho_\phi\). On the other hand, \((S_\phi^{(1)},S_\phi^{(2)})\) is a standard bivariate Gaussian pair with correlation \(\rho_\phi\). Lemma~\ref{lem:price} implies that \(L_\phi^\star\) is strictly increasing in \(\rho_\phi\), so the three maximization problems have the same solutions.
\end{proof}

\section{Why marginal stability baselines cannot give the same guarantee}

Theorem~\ref{thm:score} shows that TPR is controlled by the reliability
\(\rho_\phi\), which is a property of the joint validation law: it measures how
much of the score is reproducibly driven by the anomaly-relevant coordinate.
Marginal stability criteria only compare validation score histograms.

\begin{definition}[Marginal stability criteria]
\label{def:marginal-stability}
Let
\[
    F_{\phi,d}(t)=\Pbb(S_\phi^{(d)}\le t),
    \qquad d\in\{1,2\},
\]
be the population CDFs of the score produced by detector \(\phi\) on two
typical validation domains. A validation criterion is
\emph{marginal stability} if there exists a functional \(\Psi\) such that
\[
    M(\phi)=\Psi(F_{\phi,1},F_{\phi,2}).
\]
Equivalently, if two detectors have the same two marginal score laws, then
the criterion assigns them the same value, regardless of their paired
dependence or anomaly-detection power.
\end{definition}

The Wasserstein and threshold drift baselines are marginal stability criteria, as
their population forms are, respectively,
\[
    M_{W_1}(\phi)
    =
    W_1\!\left(
        \cL(S_\phi^{(1)}),
        \cL(S_\phi^{(2)})
    \right)
    =
    \int_0^1
    \left|
        F_{\phi,1}^{-1}(u)-F_{\phi,2}^{-1}(u)
    \right|du,
\]
and, for target FPR \(\alpha\in(0,1)\),
\[
    M_{\mathrm{thr},\alpha}(\phi)
    =
    \left|
        \Pbb\!\left(
            S_\phi^{(2)}\ge F_{\phi,1}^{-1}(1-\alpha)
        \right)-\alpha
    \right|
    =
    \left|
        1-F_{\phi,2}\!\left(F_{\phi,1}^{-1}(1-\alpha)\right)-\alpha
    \right|,
\]
assuming continuity to avoid quantile ties. In contrast, CAP does not belong to
this family because its population objective depends on the joint law of
\((S_\phi^{(1)},S_\phi^{(2)})\).

Marginal stability can measure the amount of domain shift between \(F_1\) and \(F_2\),
but it cannot determine whether the score variation responsible for anomaly
detection is reliable. This motivates the following non-identifiability result.

\begin{theorem}[Non-identifiability of TPR from marginal stability]
\label{thm:marginal-nonidentifiability}
Let \(M\) be a marginal stability criterion. Fix
\(\alpha\in(0,1)\), \(\delta>0\), and two reliability levels
\[
    0\le \rho_{\phi_-}<\rho_{\phi_+}\le 1.
\]
Let \(F_1,F_2\) be any continuous, strictly increasing validation score CDFs
with finite first moments, and let
\[
    T_d(u)=F_d^{-1}(\Phi(u)),
    \qquad d\in\{1,2\},
\]
be the monotone Gaussian maps to the target
marginal distributions. Then there exist two candidate detectors \(\phi_-\) and
\(\phi_+\) such that their validation score marginals are identical across
candidates,
\[
    F_{\phi_+,1}=F_{\phi_-,1}=F_1,
    \qquad
    F_{\phi_+,2}=F_{\phi_-,2}=F_2,
\]
but their latent reliabilities are \(\rho_{\phi_+}\) and
\(\rho_{\phi_-}\). Hence every marginal stability criterion ties the two
candidates,
\[
    M(\phi_+)=M(\phi_-),
\]
while their TPRs are strictly ordered:
\[
    \TPR_{\phi_+}(\alpha)>\TPR_{\phi_-}(\alpha).
\]
Thus marginal score stability does not identify the TPR-optimal candidate.
\end{theorem}

\begin{proof}
For each \(\rho\in[0,1]\), let
\[
    G_\rho^{(d)}
    =
    \sqrt{\rho}\,Z
    +
    \sqrt{1-\rho}\,\varepsilon^{(d)},
    \qquad d\in\{1,2\},
\]
where \(Z,\varepsilon^{(1)},\varepsilon^{(2)}\) are independent standard
Gaussian variables. Then
\((G_\rho^{(1)},G_\rho^{(2)})\) is a standard bivariate Gaussian pair with
correlation \(\rho\). Define validation scores by the monotone marginal
calibrations
\[
    S_\rho^{(d)}=T_d(G_\rho^{(d)})
    =
    F_d^{-1}(\Phi(G_\rho^{(d)})).
\]
By the probability integral transform, \(S_\rho^{(d)}\sim F_d\) for every
\(\rho\). Instantiate \(\phi_-\) with \(\rho=\rho_{\phi_-}\) and
\(\phi_+\) with \(\rho=\rho_{\phi_+}\). The two candidates have the same
marginal validation laws \(F_1,F_2\), so any marginal stability criterion
gives them the same value.

Since \(T_1\) is strictly increasing, thresholding \(S_\rho=T_1(G_\rho)\) at
FPR \(\alpha\) is equivalent to thresholding \(G_\rho\) at
\(\Phi^{-1}(1-\alpha)\). Under anomalies,
\(G_\rho\mid Y=1\sim\Normal(\delta\sqrt{\rho},1)\), so
\(\TPR_\rho(\alpha)\) is strictly increasing in \(\rho\). Therefore
\(\TPR_{\phi_+}(\alpha)>\TPR_{\phi_-}(\alpha)\).
\end{proof}

\begin{corollary}[Marginal stability can prefer the lower-power detector]
\label{cor:marginal-misranking}
Let \(\alpha\in(0,1)\), \(\delta>0\), \(\eta>0\), and candidates \(\phi_-\), \(\phi_+\) with \(0\le \rho_{\phi_-}<\rho_{\phi_+}\le 1\) so that \(\TPR_{\phi_+}(\alpha)>\TPR_{\phi_-}(\alpha)\). Let validation marginals be partially shifted
\[
    F_{\phi_-,1}=F_{\phi_-,2}=\Phi,
    \qquad
    F_{\phi_+,1}=\Phi,
    \qquad
    F_{\phi_+,2}(t)=\Phi(t-\eta).
\]
Then, \(\phi_-\) has perfect marginal stability, while \(\phi_+\) has a
location shift of magnitude \(\eta\) between validation domains. Consequently,
any marginal-stability discrepancy that is minimized by equal marginals and
increases under this location shift prefers \(\phi_-\) when smaller values are
selected. In particular, for Wasserstein distance,
\[
    M_{W_1}(\phi_-)
    =
    W_1(\Normal(0,1),\Normal(0,1))
    =
    0,
\]
whereas
\[
    M_{W_1}(\phi_+)
    =
    W_1(\Normal(0,1),\Normal(\eta,1))
    =
    \eta.
\]
Similarly, for threshold drift,
\[
    M_{\mathrm{thr},\alpha}(\phi_-)=0,
    \qquad
    M_{\mathrm{thr},\alpha}(\phi_+)
    =
    1-\Phi\!\left(\Phi^{-1}(1-\alpha)-\eta\right)-\alpha
    >
    0.
\]
\end{corollary}

\begin{proof}
Candidate \(\phi_-\) has identical validation marginals, so
both Wasserstein and threshold drift are zero. For \(\phi_+\), the two
validation marginals are \(\Normal(0,1)\) and \(\Normal(\eta,1)\). For Wasserstein,
\[
    W_1(\Normal(0,1),\Normal(\eta,1))=\eta.
\]
The threshold calibrated on the first domain is
\(\Phi^{-1}(1-\alpha)\). Under the second domain,
\[
    \Pbb(S_{\phi_+}^{(2)}\ge \Phi^{-1}(1-\alpha))
    =
    1-\Phi\!\left(\Phi^{-1}(1-\alpha)-\eta\right)
    >
    \alpha.
\]
\end{proof}

\section{Gaussian linear anomaly model}

We now instantiate the score-level result in a Gaussian mean-shift model where
the anomaly direction is analytically known. This case makes the metric
contrast explicit: CAP selects the likelihood-ratio direction when that
direction is the reproducible direction of the validation channel, while
marginal stability metrics cannot guarantee the same selection from score
histograms alone.

\begin{corollary}[Gaussian linear anomaly model]
\label{cor:linear}
Let
\[
    P_0=\Normal(0,\Sigma),
    \qquad
    P_1=\Normal(\mu,\Sigma),
    \qquad
    \Sigma\succ 0,
    \quad \mu\ne 0.
\]
Consider normalized linear scores
\[
    S_w(x)=w^\top x,
    \qquad
    w^\top\Sigma w=1,
    \qquad
    w^\top\mu\ge 0.
\]
The likelihood-ratio anomaly direction is
\[
    w_\star
    =
    \frac{\Sigma^{-1}\mu}{\sqrt{\mu^\top\Sigma^{-1}\mu}}.
\]
Let the two typical views satisfy
\[
    \begin{pmatrix}X^{(1)}\\X^{(2)}\end{pmatrix}
    \sim
    \Normal\!\left(
        0,
        \begin{pmatrix}
            \Sigma & \Gamma \\
            \Gamma^\top & \Sigma
        \end{pmatrix}
    \right),
\]
and write \(\Gamma_s=(\Gamma+\Gamma^\top)/2\). Assume that the anomaly
direction is also the unique most reproducible direction of the typical
two-view channel:
\[
    w_\star
    =
    \argmax_{w^\top\Sigma w=1,\, w^\top\mu\ge 0} w^\top\Gamma_s w.
\]
Then, for every \(\alpha\in(0,1)\),
\[
    \argmax_{w^\top\Sigma w=1,\,w^\top\mu\ge 0}
    L_w^\star
    =
    \{w_\star\}
    =
    \argmax_{w^\top\Sigma w=1,\,w^\top\mu\ge 0}
    \TPR_w(\alpha).
\]
However, any marginal stability criterion can only use the score marginals
\((F_{w,1},F_{w,2})\). Therefore, by
Theorem~\ref{thm:marginal-nonidentifiability}, such a criterion cannot
uniformly identify \(w_\star\): if a lower-power direction has the same score
marginals as \(w_\star\), the criterion must tie them, and if that direction
has smaller marginal drift, the criterion may prefer it.
\end{corollary}

\begin{proof}
For the Gaussian mean-shift model,
\[
    \log\frac{dP_1}{dP_0}(x)
    =
    \mu^\top\Sigma^{-1}x
    -
    \frac{1}{2}\mu^\top\Sigma^{-1}\mu,
\]
so the likelihood-ratio anomaly direction is \(\Sigma^{-1}\mu\). After
normalizing by \(w^\top\Sigma w=1\), this gives \(w_\star\).

For any normalized \(w\), under \(P_0\),
\[
    S_w(X)\sim\Normal(0,1),
\]
and under \(P_1\),
\[
    S_w(X)\sim\Normal(w^\top\mu,1).
\]
Therefore
\[
    \TPR_w(\alpha)
    =
    1-\Phi\!\left(\Phi^{-1}(1-\alpha)-w^\top\mu\right).
\]
Maximizing TPR is equivalent to maximizing \(w^\top\mu\) subject to
\(w^\top\Sigma w=1\). By Cauchy--Schwarz in the \(\Sigma\)-inner product,
\[
    w^\top\mu
    =
    \langle \Sigma^{1/2}w,\Sigma^{-1/2}\mu\rangle
    \le
    \sqrt{w^\top\Sigma w}\sqrt{\mu^\top\Sigma^{-1}\mu}
    =
    \sqrt{\mu^\top\Sigma^{-1}\mu},
\]
with equality only at \(w=w_\star\).

For CAP, define
\[
    Y_1=w^\top X^{(1)},
    \qquad
    Y_2=w^\top X^{(2)}.
\]
Then \((Y_1,Y_2)\) is a standard bivariate Gaussian pair with correlation
\[
    \rho(w)=\Corr(Y_1,Y_2)=w^\top\Gamma_s w,
\]
because \(w^\top\Sigma w=1\). Lemma~\ref{lem:price} implies that optimized CAP
is maximized by maximizing \(\rho(w)\). By the assumed alignment of
\(\Gamma_s\), the unique maximizer is \(w_\star\). The marginal-stability
claim is exactly the non-identifiability statement of
Theorem~\ref{thm:marginal-nonidentifiability} applied to the score marginals
of linear directions.
\end{proof}

For a concrete calculation, take \(\Sigma=I\), \(\mu=m w_\star\) with
\(\|w_\star\|=1\), and suppose the two validation domains have a benign
location shift along \(w_\star\):
\[
    X^{(1)}\sim\Normal(0,I),
    \qquad
    X^{(2)}\sim\Normal(\eta w_\star,I),
    \qquad
    \eta>0.
\]
For any admissible unit-norm direction \(w\) with \(w^\top\mu\ge 0\), write
\[
    c(w)=w^\top w_\star\in[0,1].
\]
Then
\[
    w^\top X^{(1)}\sim\Normal(0,1),
    \qquad
    w^\top X^{(2)}\sim\Normal(\eta c(w),1).
\]
The anomaly power is
\[
    \TPR_w(\alpha)
    =
    \Phi\!\left(m c(w)-\Phi^{-1}(1-\alpha)\right),
\]
so \(w_\star\) is optimal because \(c(w_\star)=1\). But the marginal stability
baselines are minimized away from \(w_\star\). Specifically,
\[
    M_{W_1}(w)=\eta c(w),
\]
and
\[
    M_{\mathrm{thr},\alpha}(w)
    =
    1-\Phi\!\left(\Phi^{-1}(1-\alpha)-\eta c(w)\right)-\alpha,
\]
which are both strictly increasing in \(c(w)\) on \([0,1]\). Hence, when the
dimension is at least two, every direction in the region
\[
    \mathcal R_\epsilon
    =
    \{w:\|w\|=1,\ 0\le c(w)\le 1-\epsilon\},
    \qquad
    \epsilon\in(0,1),
\]
has lower Wasserstein and threshold drift than \(w_\star\), even though every
such direction has strictly lower TPR than \(w_\star\). This is the concrete
Gaussian case in which marginal stability prefers a region of lower-power
scores over the true anomaly direction.

\end{document}
